\chapter{Sorting and search}
\MARSChapterHeading{Sorting and search}
\label{revision-lab:algorithms-and-cost}

Read the main chapter, then import the student package
\texttt{009-sorting-costs} and continue its downloaded lesson in Learn. Answer
the inline predictions using their source pointers, then use the imported
\texttt{sorting-costs.asm} for the guided Debug sequence. Keep its default
insertion algorithm and eight-word count for those observations. Use separate
copies for the smaller inputs, selection sort and deliberate faults below.

The worked source is also available at
\path{revision/labs/files/sorting-costs.asm} for ordinary Run and Debug. Read
the two invariants before following the loops. The purpose is to explain a
measured difference, including the work done to show the algorithm. The native
lesson checks register observations; this lab connects them with complete array,
metadata and framebuffer state and the public measurement tool.

\section{Predict complete exchanges}

In a copy of the program, set \texttt{count} to four and retain the first four
words $[2,-2,3,-4]$. Set \texttt{algorithm} to zero for insertion. Write the
array after each complete exchange, and mark the boundary of the sorted prefix
after each outer iteration. An outer iteration can contain several exchanges
or none. Compare your trace with the main chapter's table only after deriving it.

Repeat with \texttt{algorithm} set to one for selection. At each outer iteration,
record the first index of the suffix's minimum before making the exchange. Explain
why both methods produce the same four active words while leaving the other four
words untouched. If your result is sorted, also check that every original value
occurs the same number of times. Sorted order by itself cannot detect lost values.

Restore the default eight-word count. Predict which method will make more exchanges
on this input. Then predict the count for an already sorted copy and a reversed
copy of the same eight distinct keys. Keep the input values fixed while changing
their order: changing the bar heights at the same time would add another cause
for a change in rendering work.

\section{Inspect a half-exchange and a publication}

Use Debug with insertion selected. Break at \texttt{insertion\_compare}. On its
first visit, identify the right and left keys in \texttt{\$t1} and
\texttt{\$t2}, and derive both store addresses from \texttt{\$t0}. The right key
is $-2$ and the left key is two. The signed comparison permits the exchange.

The guided sequence begins at the initial \texttt{sort\_ready}, then reaches
this first comparison. Its next stop is the second \texttt{sw} under
\texttt{insertion\_compare\_end}, before that store executes. Move your single
line breakpoint to that instruction and resume. To reach the same boundary
manually, step through the branch and the following \texttt{nop}, then execute
just the first store. Inspect the array: its first two words now both contain $-2$.
The value two still exists in a register, awaiting the second store. Explain
why this temporary memory image is not a completed sorting result. Execute the
second store and stop at \texttt{insertion\_exchanged}. The prefix now starts
$[-2,2,3,-4]$, but phase, exchange count and framebuffer still describe the initial
publication. Those fields are updated by the helper call that follows.

Figure~\ref{fig:revision-sort-half-exchange} shows the actual pause before the
second store. The hexadecimal value in \texttt{\$t1} is the sign-extended $-2$;
\texttt{\$t2} still contains two. The memory inspection and the register
inspection answer different questions: what is already stored, and what remains
available to finish the operation.

\begin{figure}[htbp]
\centering
\begin{tikzpicture}[alt={Two details from one light-theme Debug photograph before the second store of insertion's first exchange. The first store has completed. The pending instruction is sw t2,0(t0). Register t1 contains hexadecimal fffffffffffffffe, signed minus two; t2 contains hexadecimal 0000000000000002, signed two. Outlines identify the pending instruction and the retained key. The temporary array duplicate is explained in the surrounding text, not displayed in these crops.}]
\node[anchor=north west,inner sep=0,font=\small\bfseries] (codetitle) at (0,0) {1. The second store is pending};
\node[anchor=north west,inner sep=0,font=\small\bfseries] (regtitle) at ([xshift=0.465\linewidth]codetitle.north west) {2. The displaced key is retained};
\node[anchor=north west,inner sep=0] (code) at ([yshift=-2mm]codetitle.south west) {\MARSRevisionCrop{revision/captures/sorting/native-half-exchange.png}{0.43\linewidth}{0.426562500}{0.501111111}{0.364062500}{0.428888889}};
\node[anchor=north west,inner sep=0] (regs) at ([yshift=-2mm]regtitle.south west) {\MARSRevisionCrop{revision/captures/sorting/native-half-exchange.png}{0.50\linewidth}{0.225000000}{0.031111111}{0.525000000}{0.897777778}};
\begin{scope}[shift={(code.south west)},x={($(code.south east)-(code.south west)$)},y={($(code.north west)-(code.south west)$)}]
\draw[white,line width=3pt] (0.34,0.02) rectangle (0.97,0.40);
\draw[black,line width=0.9pt] (0.34,0.02) rectangle (0.97,0.40);
\end{scope}
\begin{scope}[shift={(regs.south west)},x={($(regs.south east)-(regs.south west)$)},y={($(regs.north west)-(regs.south west)$)}]
\draw[white,line width=3pt] (0.002,0.02) rectangle (0.99,0.35);
\draw[black,line width=0.9pt] (0.002,0.02) rectangle (0.99,0.35);
\end{scope}
\end{tikzpicture}
\caption{An automated Debug capture between the two stores. Both first array words already contain $-2$, but \texttt{\$t2} still holds two for the pending second store. The register panel shows full-width hexadecimal values. Outlines mark the pending store and the retained key.}
\label{fig:revision-sort-half-exchange}
\end{figure}


Continue to \texttt{sort\_ready}. Phase should be one, the count should be one,
and the frame should describe the exchanged array. Derive one white and one
black pixel from $y\geq3-values[x]$ instead of judging the shape alone. The first
visit to this label, before sorting, has phase zero and count zero; distinguish
it from this first exchange visit when using a breakpoint from program entry.

Follow the return path through \texttt{render} and \texttt{publish\_exchange}.
The renderer may change argument and temporary registers. Find the insertion
cursor and array pointer that survive in saved registers, and inspect the saved
return links in the nested frames. Continuing the loop must not depend on a
temporary register retaining a value across either call.

The next guided stop is \texttt{insertion\_outer} at index two. Its first two
words form the sorted prefix, while the next key is about to be processed. The
inner cursor reached zero on the previous iteration. The publication helper has
returned, so only the sorter's 48-byte frame remains.

After the sorter returns, stop at \texttt{sort\_end}. The caller's stack pointer
is restored and all eight words are sorted, but phase still describes the last
exchange publication. The final guided stop, \texttt{done}, follows the caller's
phase-two publication. Finish or stop Debug before ordinary Run. Learn retains
predictions and checks across project closure; edits require fresh execution
evidence for the changed source.

At the final visit, phase is two and the count is 17 for insertion or five for
selection on the default input. Both final arrays are
$[-4,-3,-2,-1,0,1,2,3]$. Count the initial and final frames as well as exchange
frames. A fast Run need not repaint the tool after every render, so count metadata
or traced stores rather than visible animation frames.

\section{Explain a boundary or ordering fault}

In a disposable copy, replace insertion's signed key comparison with an unsigned
comparison. Predict what happens when a positive key is compared with a negative
key. Use the loaded 64-bit register patterns, remembering that \texttt{lw} sign
extends its word. Restore the signed comparison and show where the first incorrect
exchange decision disappears.

Next inspect the check for a zero insertion cursor. State the invalid address that
the left-neighbor load would use if that check were removed. Test lengths zero and
one before trying a long input. Neither needs a key comparison, and neither gives
the sorter permission to read the word before the array or the inactive suffix.
A readable address is not necessarily part of the supplied input.

For selection, use repeated minimum values. Explain why the current comparison
retains the first minimum rather than the last equal one. This policy fixes the
exchange trace, but does not make long-distance selection exchanges stable. Trace
$[2_a,2_b,1]$ using the subscripts as conceptual record identities. The two equal
keys can change relative order even though their integer values look identical.

Finally, skip a black pixel store in a copy of \texttt{render}. Arrange for a bar
to shrink between publications. Compare the array with the frame and locate the
stale pixel. The array can be sorted correctly while its observation is wrong.
Restore the store and confirm that rendering overwrites all 64 cells, including
inactive columns.

\section{Measure the work that answers your question}

From the public repository root, run the supplied measurement tool with JDK 25
and the MARS CLI archive:
\begin{verbatim}
python3 tools/measure_revision_sorting.py \
  --java /path/to/jdk-25/bin/java \
  --cli /path/to/cli.jar --output /new/report/directory
\end{verbatim}
The output directory must be new so that one run cannot silently replace an earlier
receipt. Retain the JSON report together with its source and CLI hashes. The tool
runs both algorithms on seven input patterns, including empty and singleton
prefixes. It checks output state while counting completed simulator instructions,
key comparisons and framebuffer stores.

Read the mixed-input rows first. Explain how insertion can use 344 instructions
inside its routine while selection uses 399, yet the complete visualized runs use
19,633 and 7,510 respectively. Derive the 19 versus seven render calls from the
exchange counts, then multiply by 64 to obtain 1,216 versus 448 framebuffer stores.
Each exchange adds observation work outside the measured sorting-routine interval.

Write a one-sentence scope statement for each number you report. The routine
interval includes setup, loop control, argument preparation, call instructions and
the return. It excludes the helper and renderer bodies. The whole-run interval
includes those bodies, initial and final publications, and exit. Neither interval
measures assembly, session startup, elapsed host time or hardware cycles.

Use the sorted and reverse-ordered rows to distinguish input sensitivity from
scope. Selection makes 28 key comparisons on each eight-word input; insertion
makes seven on the sorted input and 28 on the reversed input. On an empty prefix,
there are no exchanges or key comparisons, but the program still draws its initial
and final frames. Explain why a zero-length sort can therefore have a nonzero
whole-run instruction count without performing an invalid array read.

If you change an algorithm, first establish that its result and interface remain
correct. Then repeat the same inputs and scope. Changing the renderer, publication
policy and sorter simultaneously makes a lower total harder to attribute. A
performance claim should identify which behavior changed and which work was counted.

\section{Write a stable sign partition independently}

Open the imported \texttt{independent/partition-negative.asm}. Implement
\texttt{partition\_negative}
with an array pointer in \texttt{\$a0} and a length from zero through eight in
\texttt{\$a1}. Return the number of negative words in \texttt{\$v0}. Place the
negative words first and nonnegative words second, retaining their original
relative order within each group. Zero is nonnegative. Unlike the visualizer,
this task accepts every signed 32-bit word value.

For the default input, write the expected result before coding:
$[-2,-4,-3,-1,2,3,0,1]$, with a returned count of four. Notice that the negative
group is not numerically sorted. An ordinary ascending sort would put all negative
words first, but would violate this task's stability requirement. The condition
you must preserve comes from the interface, not from the previous worked algorithm.

Choose a method and state its invariant. For an in-place approach, consider where
to insert a newly encountered negative word so that earlier negative words and
nonnegative words retain their order. For a temporary-buffer approach, account
for every input word exactly once in each required pass, and restore the stack
pointer after copying back. Either approach is acceptable if it meets the contract.

Preserve saved registers, the stack pointer, caller-owned stack bytes, guards and
inactive words. Limit array loads and stores to the supplied span. The pointer
can name an offset within the allocation, and the supplied length is authoritative;
do not replace either argument with the global test fixture. Private stack storage
is permitted, but writes into the caller's frame are not.

Test empty input, one negative word, one zero, all-negative and all-nonnegative
inputs, repeated values and the most negative signed word. Explain why comparing
with zero directly avoids needing to negate that extreme value. Use \textbf{Check this program} in Learn to check the current editor text,
including an unsaved edit. The unfinished starter should fail. For the same
public cases outside the IDE, run from the repository root:
\begin{verbatim}
java -jar /path/to/cli.jar grade \
  revision/lessons/009-sorting-costs/independent/partition-negative.kts \
  /path/to/your/partition-negative.asm
\end{verbatim}
The 56 cases include all 45 valid offset/length pairs in an eight-word fixture,
alongside signed extremes and repeated keys. They pass a length different from
the preserved global \texttt{count} word. A pointer to the middle of the array
and its passed length define the active span; the fixture's base and global
count do not replace that interface.

An explicit guard precedes the imported array. Access histories watch that
guard, the complete backing array and its following guard, requiring overlapping
reads and stores to fit wholly within the active span. Separate store histories
protect metadata, guards and the 32 caller-owned stack bytes above the aligned
incoming pointer. A same-value or later-restored caller store is still a
violation. Private frames below that pointer remain permitted. These bounded
observations do not prescribe a frame layout or claim to audit every address.

Each case calls through both supplied continuations. The second call receives
the already partitioned span, which should retain its values and negative
count. A pass establishes 112 returned calls on the public inputs, rather than
every possible signed array. Keep an invariant and counterexample beside the
report.

The starter deliberately stops until you replace its missing function. On a failed
case, identify whether the disagreement concerns group order, relative order,
returned count, access bounds or preserved state. Repair that requirement, then
rerun the whole check set. A correct default picture or one correct output array
does not establish the other parts of the contract.

\section{Companion lab: a shortest route through a maze}
\label{revision-lab:maze}

Sorting required a useful order for array entries. A search also needs an order
for unfinished work. In this lab, an array of cell IDs records places waiting to
be explored. Taking the oldest waiting cell first lets the search expand outward
in layers of increasing distance. The result is a shortest route through an
unweighted maze, together with enough intermediate state to explain that route.

This is original project \texttt{010-maze-explorer}. Its contract is
\texttt{assignment.md} in that package. The revised work area is
\texttt{maze-search.asm}, under \texttt{revision/labs/files}. Both public checks
remain in the original package's \texttt{checks} directory. Implement the search and its
\texttt{render} function. The starter's \texttt{break} instructions deliberately
stop execution until those missing bodies are replaced. Organize private functions
and registers as you choose; the public allocations, observations and renderer
interface establish what another program can inspect.

You will use byte arrays, row-major addresses, loop invariants and calls that
preserve live state. The new ideas are graph adjacency, discovery-time marking,
a first-in, first-out frontier, predecessor links and a shortest-distance argument.
Breadth-first search, or BFS, is the conventional name for this layer-by-layer
search.\footnote{For the layer argument and predecessor reconstruction, see MIT
OpenCourseWare, \emph{6.006 Introduction to Algorithms}, Spring 2020, Recitation 9,
pp.~3--4: \url{https://ocw.mit.edu/courses/6-006-introduction-to-algorithms-spring-2020/resources/mit6_006s20_r09/}.}

\subsection*{Represent cells and legal moves separately}

The input is an eight-by-eight array of bytes. A zero \texttt{walls} byte denotes
an open cell and one denotes a wall. The word inputs \texttt{start} and
\texttt{goal} name open cells by IDs from zero through 63. A cell at column $x$
and row $y$ has ID $8y+x$. Its byte entry is at $walls+ID$; its pixel word is at
$frame+4ID$. An ID is neither of those addresses.

Treat each open cell as a vertex of a graph, and each legal one-cell move as an
edge. Moves are north, east, south or west, with cost one each. Diagonals are
excluded. The graph is implicit: the program computes possible neighbors from
coordinates rather than storing a linked collection for every cell.

Derive $x=ID\bmod8$ and $y=\lfloor ID/8\rfloor$, apply a direction displacement,
then reject a coordinate outside zero through seven. Only after these checks
should the candidate be flattened back to an ID. Cell seven's eastern candidate
has column eight and must be rejected. Adding one to its ID gives eight, which
is a valid array index but belongs to the next row. An array bound alone cannot
establish adjacency.

The complete \texttt{maze-neighbor.asm} experiment isolates this geometry. Its
\texttt{neighbor\_id} function takes a valid cell ID in \texttt{\$a0} and a
direction index in \texttt{\$a1}: zero north, one east, two south, three west.
It returns the neighbor ID in \texttt{\$v0}, or $-1$ when the move leaves the grid.
It does not inspect walls; an in-bounds cell can still be blocked. The direction
tables contain signed bytes, so \texttt{lb} must preserve the $-1$ displacements.

\noindent\begin{minipage}{\linewidth}
\MARSCodeLabel{Check both coordinates before producing an ID}
\begin{minted}{mars-idea}
neighbor_id:
    andi $t0, $a0, 7
    srl $t1, $a0, 3
    la $t2, dx
    daddu $t2, $t2, $a1
    lb $t3, 0($t2)
    addu $t0, $t0, $t3
    bltzc $t0, neighbor_absent
    nop
    li $t4, 8
    bgec $t0, $t4, neighbor_absent
    nop
    la $t2, dy
    daddu $t2, $t2, $a1
    lb $t3, 0($t2)
    addu $t1, $t1, $t3
    bltzc $t1, neighbor_absent
    nop
    bgec $t1, $t4, neighbor_absent
    nop
    sll $v0, $t1, 3
    addu $v0, $v0, $t0
    jrc $ra
neighbor_absent:
    li $v0, -1
    jrc $ra
\end{minted}
\end{minipage}

The bit mask extracts the low three column bits, and the shift extracts the row.
After applying a displacement, the two signed checks on each coordinate reject
negative and too-large results. Pointer additions use full-width instructions;
the coordinates themselves fit comfortably in signed words. Saved registers and
the stack pointer are untouched. This helper is an experiment you may call or
adapt, not a required private decomposition of the full search.

Run \texttt{revision/labs/checks/maze-neighbor.kts} against that complete source
with the CLI. It checks all 256 valid cell/direction combinations and the supplied
caller. Before running it, predict east from seven, west from eight, north from
zero, and south from 55. Explain each answer with coordinates. In a copy, replace
\texttt{lb} with \texttt{lbu} and find which directions now fail. Restore the
signed loads before using the experiment as a search component.

\subsection*{Make discovery a single event}

Five 64-byte arrays hold different facts. \texttt{states} contains zero for unseen,
one for discovered and waiting, and two for processed. \texttt{parents} stores the
ID of the cell that first discovered an entry. \texttt{distance} stores its number
of edges from the start. \texttt{queue} retains discovery IDs in FIFO order, and
\texttt{path} eventually marks the chosen route. Unknown parents and distances
are 255; queue and path bytes begin at zero. Initialize every entry on every launch,
including those that the eventual search will never reach.

These arrays are byte-indexed, so load IDs and distances with \texttt{lbu}.
The sentinel 255 is not a valid ID and is not a distance to add one to. A simple
route in a graph of at most 64 cells has at most 63 edges, leaving the sentinel
outside the valid distance range. Word metadata uses \texttt{lw}/\texttt{sw};
all public byte-array bases and metadata words are four-byte aligned. Preserve
the declared guard words and execute no stores to \texttt{walls}, \texttt{start}
or \texttt{goal}, even if later writes could restore them.

Initialize the start as waiting with distance zero and no parent, then append
its ID. For each newly discovered neighbor, set its waiting state, parent and
distance before appending it. The parent is the current cell; the distance is
that current cell's known distance plus one. Skip walls and every already
discovered cell, including waiting cells that have not yet been removed.

The four-cell room shows why the mark belongs at discovery. Zero discovers both
one and eight.
Processing one can discover nine. When eight later examines nine, it must see the
existing mark and leave nine's parent and queue entry alone. If marking waited
until dequeue, both predecessors could append nine before either copy was
processed. Cycles would repeat work, invalidate unique history and undermine the
queue capacity argument.

\subsection*{Use the queue to explain the shortest distance}

The monotone queue uses two word indices. \texttt{head} names its next unread
byte and \texttt{tail} its first unused byte. Pending work occupies the half-open
interval $[head,tail)$; retained history occupies $[0,tail)$. The invariant is
$0\leq head\leq tail\leq64$. An empty pending queue has equal indices, even
though earlier queue entries still contain IDs. Do not clear a removed entry.
The unused suffix remains zero.

This queue does not wrap like the preceding ring assignment. Each cell is
appended at most once, so the entire search needs at most 64 slots, regardless
of how many slots remain pending at one moment. Appending ID zero is valid:
use head and tail to identify occupied entries instead of treating a zero byte
as an end marker.

The FIFO order is what gives the distance its meaning. Initially only distance
zero is waiting. When a cell at distance $d$ is removed, its unseen neighbors
receive $d+1$ and go behind work already waiting. All cells at smaller distances
are therefore processed before any newly appended farther layer. If a neighbor
had a shorter route, the preceding cell on that route would have been processed
in an earlier layer and discovered it already. Its first assigned distance is
thus shortest. This argument depends on equal edge costs and FIFO removal;
a stack or weighted terrain would require a different argument.

\begin{figure}[htbp]
  \centering
  \begin{tikzpicture}[x=17mm,y=-17mm,alt={The upper-left two by two room contains cell zero at top left, one at top right, eight at bottom left and nine at bottom right. All other maze cells are walls. Orthogonal links form a cycle. Two routes from zero to nine have two edges: through one or through eight. Arrows highlight the first route; a parent walk traverses it backward from nine to one to zero.}]
    \draw[thick] (0,0) rectangle (2,2);
    \draw (1,0)--(1,2) (0,1)--(2,1);
    \node at (.5,.5) {0}; \node at (1.5,.5) {1};
    \node at (.5,1.5) {8}; \node at (1.5,1.5) {9};
    \draw[->,very thick] (.72,.5)--(1.28,.5);
    \draw[->,very thick] (1.5,.72)--(1.5,1.28);
    \node[anchor=west,align=left,text width=62mm] at (2.4,.55) {Two equal-length routes:\\$0\to1\to9$\\$0\to8\to9$};
    \node[anchor=west,align=left,text width=62mm] at (2.4,1.5) {For the highlighted route,\\parents lead backward:\\$9\to1\to0$.};
  \end{tikzpicture}
  \caption{Cell IDs use the full maze's eight-cell row stride, even in this small open room. A route chooses edges; a parent entry records the preceding cell's ID.}
  \label{fig:revision-maze-room}
\end{figure}

\begin{table}[htbp]
  \centering
  \small
  \begin{tabular}{lrrlll}
    \toprule
    Boundary & Head & Tail & Queue history    & Waiting & New parent           \\
    \midrule
    Initial  & 0    & 1    & \texttt{0}       & 0       & none                 \\
    After 0  & 1    & 3    & \texttt{0,1,8}   & 1,8     & $1\gets0$, $8\gets0$ \\
    After 1  & 2    & 4    & \texttt{0,1,8,9} & 8,9     & $9\gets1$            \\
    After 8  & 3    & 4    & \texttt{0,1,8,9} & 9       & none                 \\
    After 9  & 4    & 4    & \texttt{0,1,8,9} & none    & none                 \\
    \bottomrule
  \end{tabular}
  \caption{A FIFO trace in the four-cell room, using north, east, south, west neighbor order. Head names the next unread queue position; tail is the first unused position. History is retained after removal. The final path is $0,1,9$, with two edges.}
  \label{tab:revision-maze-frontier}
\end{table}


Follow Table~\ref{tab:revision-maze-frontier} by hand before implementing the
loop. With north, east, south, west order, nine first receives parent one.
With west, south, east, north order, the waiting order after zero is eight, one,
and nine first receives parent eight. Both routes have two edges. The public
checks accept either order and follow your actual queue history; they require a
shortest valid chain, not the reference program's tie choice.

Remove one pending ID, advance head, and mark that current cell processed.
If it is the goal, skip its neighbor expansion. Otherwise inspect all four
candidates, appending only unseen open neighbors. Publish after the complete
expansion, then choose the next item. A reachable goal is finalized when dequeued
under this contract, even if its shortest distance became known earlier.

\subsection*{Publish a state that another observer can interpret}

Set phase zero after initialization, phase one after each complete dequeue and
expansion, and phase two after success or exhaustion. At every publication,
\texttt{head}, \texttt{tail} and \texttt{current} describe that boundary. Initially
current is $-1$. Result and length remain $-1$ throughout phases zero and one,
including the phase-one publication of a dequeued goal. The path mask remains
zero until reconstruction. A cell in the middle of expansion can have processed
state internally, but phase-one observers must receive all its discoveries together.

The leaf-style \texttt{render} interface reads \texttt{walls}, \texttt{states}
and \texttt{path}, writes all 64 pixel words at \texttt{frame}, and changes no
input array or metadata. It preserves full saved-register values and the stack
pointer. It reaches \texttt{maze\_ready} exactly once with the frame complete
and preserved state restored, then returns through the supplied return address.
The \texttt{done} label supports an independent call. If you introduce nested
calls or temporary stack storage, apply the earlier call-frame rules.

Color priority is wall, path, then state. In hexadecimal RGB, walls are
\texttt{202020}, path cells \texttt{0072B2}, unseen cells \texttt{FFFFFF}, waiting
cells \texttt{E69F00}, and processed cells \texttt{009E73}. Repaint every pixel,
including pixels whose new color is white or dark. The colors help locate state;
array values and IDs provide the explanation without depending on color perception.
The framebuffer is the first static allocation at hexadecimal \texttt{10010000};
its immediate guard is at \texttt{10010100}. Other named arrays may move.

Set a breakpoint at \texttt{maze\_ready} in Debug. Inspect phase zero and the
single waiting start, then continue through the four-cell trace. At each phase-one
pause, derive the pending range from head and tail and inspect the newly appended
IDs and their parents. Step beyond the current occurrence before waiting for the
next one. A fast Run may show only the final frame; screen refreshes are not a
count of completed expansions.

\subsection*{Follow parents and distinguish failure from a zero-length route}

On success, set result to one and length to the goal's distance. Starting at the
goal, mark that cell in \texttt{path}, then follow its parent until reaching the
start. Mark the start too, and stop before reading its sentinel parent. Each
parent must be an adjacent open cell whose distance is one smaller. This decreasing
nonnegative distance explains why the chain terminates without a cycle.

For route $0,1,9$, reconstruct backward as $9,1,0$. The mask has three marked
cells but length is two: the contract counts edges. If start equals goal, the
search still initializes and dequeues that cell, then succeeds with length zero
and one marked cell. It does not expand the goal's neighbors.

If pending work becomes empty without removing the goal, set result to zero,
retain length $-1$ and leave the path mask all zero. Reach phase two, redraw and
exit normally in either case. An unreachable goal and a zero-edge success have
different result, length and path states; do not infer either from an empty queue
alone.

Figure~\ref{fig:revision-maze-routes} compares the two valid reconstructed chains.
Use the IDs and publication outlines as well as the colors when comparing your
result. A different parent for nine is acceptable when its chain remains valid
and shortest.

\begin{figure}[htbp]
\centering
\begin{tikzpicture}[alt={Two actual native Bitmap Display observations cropped to the upper-left two by two open room. The left route marks cells zero, one and nine; cell eight is processed but off the route. The right route marks zero, eight and nine; cell one is processed but off the route. Native publication numbers identify all four cells and double outlines identify the three route cells in each panel, independently of color. Both routes have two edges.}]
\node[anchor=north west,inner sep=0,font=\small\bfseries] (lefttitle) at (0,0) {1. East before south};
\node[anchor=north west,inner sep=0] (left) at ([yshift=-2mm]lefttitle.south west) {\MARSRevisionCrop{revision/captures/maze/maze-east-first.png}{0.28\linewidth}{0.203906250}{0.698333333}{0.648046875}{0.091111111}};
\coordinate (rightorigin) at ([xshift=12mm]left.north east);
\node[anchor=north west,inner sep=0,font=\small\bfseries] (righttitle) at (rightorigin |- lefttitle.north west) {2. South before east};
\node[anchor=north west,inner sep=0] (right) at (righttitle.south west |- left.north west) {\MARSRevisionCrop{revision/captures/maze/maze-south-first.png}{0.28\linewidth}{0.203906250}{0.698333333}{0.648046875}{0.091111111}};
\foreach \panel in {left,right} {
 \begin{scope}[shift={(\panel.south west)},x={($(\panel.south east)-(\panel.south west)$)},y={($(\panel.north west)-(\panel.south west)$)}]
 \draw[black,line width=.4pt] (0,0) rectangle (1,1) (.5,0)--(.5,1) (0,.5)--(1,.5);
 \foreach \x/\y/\id in {.25/.75/0,.75/.75/1,.25/.25/8,.75/.25/9} {\node[font=\normalsize,fill=white,inner sep=1.5pt] at (\x,\y) {\id};}
 \end{scope}
}
\foreach \panel/\x/\y in {left/.02/.52,left/.52/.52,left/.52/.02,right/.02/.52,right/.02/.02,right/.52/.02} {
 \begin{scope}[shift={(\panel.south west)},x={($(\panel.south east)-(\panel.south west)$)},y={($(\panel.north west)-(\panel.south west)$)}]
 \draw[black,line width=1.8pt] (\x,\y) rectangle ++(.46,.46);
 \draw[white,line width=.7pt] (\x,\y) rectangle ++(.46,.46);
 \end{scope}
}
\end{tikzpicture}
\caption{Actual light-theme Debug observations of the open room; other cells are walls. Native LaTeX numbers identify cells, and double outlines mark each route. Neighbor order changes the chosen shortest chain; both edge counts are two.}
\label{fig:revision-maze-routes}
\end{figure}


In a copy, make only cells seven and eight open, with seven the start and eight
the goal. A result of one reveals row wrapping. Use cells zero and nine for a
diagonal-only case; they are disconnected too. For the original four-cell room,
delay the discovery mark in a disposable implementation and inspect the first
duplicate append. Repair that transition before trying a larger open maze.

\subsection*{Build and check the search independently}

Begin with full initialization and the phase-zero frame. Then implement one
dequeue and its complete expansion; use the table to inspect the first two such
boundaries. Generalize that loop, add exhaustion and success, and finally reconstruct
the predecessor chain. These stages separate faulty geometry, duplicate discovery
and incorrect reconstruction instead of hiding all three behind one final picture.
Keep the renderer independently callable so its faults can be investigated without
running a complete search.

Run both package checks against your revised source:
\begin{verbatim}
java -jar /path/to/cli.jar grade \
  packages/010-maze-explorer/checks/public.kts \
  /path/to/your/maze-search.asm
java -jar /path/to/cli.jar grade \
  packages/010-maze-explorer/checks/render.kts \
  /path/to/your/maze-search.asm
\end{verbatim}
The checks inspect every published transition and accept your scratch-register
choices. Renderer checks supply independent arrays and saved registers. Keep a
failing input and locate its first disagreement: a missing neighbor, duplicate
append or changed parent can precede the wrong final mask. Repair that transition
and rerun both checks, rather than patching only the picture.

After the basic search passes, change the open room to force a detour and justify
its new edge count. Predict which parent changes if you reverse neighbor order.
Then measure the largest pending queue size at completed expansion boundaries;
state that sampling boundary, since it excludes the internal moment after a
removal but before its new appends. For a larger extension, design a rectangular
interface with explicit row stride and capacity. Do not carry the eight-column
mask or 64-slot proof into that interface without deriving their replacements.

